\documentclass[unicode, 12pt, a4paper,oneside,fleqn]{article}

\usepackage{../src/styles}
\begin{document}

\begin{titlepage}
\begin{center}
\bfseries

{\Large Московский авиационный институт\\ (национальный исследовательский университет)

}

\vspace{48pt}

{\large Факультет информационных технологий и прикладной математики
}

\vspace{36pt}

{\large Кафедра вычислительной математики и~программирования

}


\vspace{48pt}

Лабораторная работа \textnumero 1 по курсу \enquote{Дискретный анализ}

\end{center}

\vspace{72pt}

\begin{flushright}
\begin{tabular}{rl}
Студент: & И.\,А. Казаков \\
Преподаватель: & А.\,А. Кухтичев \\
Группа: & М8О-206БВ-24 \\
Дата: & \\
Оценка: & \\
Подпись: & \\
\end{tabular}
\end{flushright}

\vfill

\begin{center}
\bfseries
Москва, \the\year
\end{center}
\end{titlepage}

\pagebreak

\CWHeader{Лабораторная работа \textnumero 1}

\CWProblem{
Требуется разработать программу, осуществляющую ввод пар \enquote{ключ-значение}, их 
упорядочивание по возрастанию ключа указанным алгоритмом сортировки за линейное время и вывод отсортированной последовательности.

{\bfseries Вариант сортировки:} Карманная сортировка.

{\bfseries Вариант ключа:} { \normalfont\ttfamily Вещественные числа в диапазоне $[-100, 100]$. }

{\bfseries Вариант значения:} { \normalfont\ttfamily Строки переменной длины до 2048 символов.}
}
\pagebreak

\section{Описание}
Требуется написать реализацию алгоритма карманной сортировки.

Как сказано в \cite{Kormen}: \enquote{идея, лежащая в основе карманной сортировки, заключается в том, чтобы разбить интервал на n одинаковых интервалов, а затем
распределить по этим карманам n входных величин $x$}.

\pagebreak

\section{Исходный код}
На каждой непустой строке входного файла располагается пара \enquote{ключ-значение}, поэтому создадим  
класс \texttt{TPair}, в которой будем хранить вещественный ключ и значение. В начале работы все элементы 
считываются в динамический массив, после чего выполняется подсчет количества элементов для каждого кармана 
и вычисление их индекса в итоговом массиве. Затем данные распределяются по карманам и сортируются через сортировку вставкой.

\begin{lstlisting}[language=C++]
#include <algorithm>
#include <chrono>
#include <iomanip>
#include <iostream>
#include <string>

const double MIN_KEY = -100.0;
const double MAX_KEY = 100.0;
const double RANGE = MAX_KEY - MIN_KEY;

class TPair {
   public:
    double key;
    std::string value;

    TPair() : key(0.0), value("") {}

    void Print() const {
        std::cout << std::fixed << std::setprecision(6) << key << '\t' << value << "\n";
    }
};

void ArrayInsertionSort(TPair* arr, int n) {
    for (int i = 1; i < n; ++i) {
        TPair key = std::move(arr[i]);
        int j = i - 1;
        while (j >= 0 && arr[j].key > key.key) {
            arr[j + 1] = std::move(arr[j]);
            j--;
        }
        arr[j + 1] = std::move(key);
    }
}

int main() {
    std::ios_base::sync_with_stdio(false);
    std::cin.tie(nullptr);

    int startSize = 1024;
    int realSize = 0;
    TPair* raw = new TPair[startSize];

    double inputKey;
    std::string inputValue;

    while (std::cin >> inputKey) {
        std::cin.ignore(1, '\t');
        if (!std::getline(std::cin, inputValue)) {
            break;
        }

        if (realSize >= startSize) {
            startSize *= 2;
            TPair* temp = new TPair[startSize];
            for (int i = 0; i < realSize; ++i) {
                temp[i] = std::move(raw[i]);
            }
            delete[] raw;
            raw = temp;
        }
        raw[realSize].key = inputKey;
        raw[realSize].value = std::move(inputValue);
        realSize++;
    }

    if (realSize == 0) {
        delete[] raw;
        return 0;
    }

    TPair* stlCopy = new TPair[realSize];
    for (int i = 0; i < realSize; ++i) {
        stlCopy[i] = raw[i];
    }

    std::chrono::high_resolution_clock::time_point startBucket = std::chrono::high_resolution_clock::now();

    int bucketCount = realSize;
    int* counts = new int[bucketCount]();

    for (int i = 0; i < realSize; ++i) {
        int idx = (int)((raw[i].key - MIN_KEY) / RANGE * (bucketCount - 1));
        if (idx < 0) {
            idx = 0;
        }
        if (idx >= bucketCount) {
            idx = bucketCount - 1;
        }
        counts[idx]++;
    }

    int* offsets = new int[bucketCount];
    offsets[0] = 0;
    for (int i = 1; i < bucketCount; ++i) {
        offsets[i] = offsets[i - 1] + counts[i - 1];
    }

    TPair* sorted = new TPair[realSize];
    int* currentOffset = new int[bucketCount];
    for (int i = 0; i < bucketCount; ++i) {
        currentOffset[i] = offsets[i];
    }

    for (int i = 0; i < realSize; ++i) {
        int idx = (int)((raw[i].key - MIN_KEY) / RANGE * (bucketCount - 1));
        if (idx < 0) {
            idx = 0;
        }
        if (idx >= bucketCount) {
            idx = bucketCount - 1;
        }
        sorted[currentOffset[idx]++] = std::move(raw[i]);
    }

    for (int i = 0; i < bucketCount; ++i) {
        if (counts[i] > 1) {
            ArrayInsertionSort(sorted + offsets[i], counts[i]);
        }
    }

    std::chrono::high_resolution_clock::time_point endBucket = std::chrono::high_resolution_clock::now();

    std::chrono::high_resolution_clock::time_point startStl = std::chrono::high_resolution_clock::now();
    std::stable_sort(stlCopy, stlCopy + realSize, [](const TPair& a, const TPair& b) {
        return a.key < b.key;
    });
    std::chrono::high_resolution_clock::time_point endStl = std::chrono::high_resolution_clock::now();

    std::chrono::duration<double, std::milli> bucketTime = endBucket - startBucket;
    std::chrono::duration<double, std::milli> stlTime = endStl - startStl;

    bool match = true;
    for (int i = 0; i < realSize; ++i) {
        if (sorted[i].key != stlCopy[i].key || sorted[i].value != stlCopy[i].value) {
            match = false;
            break;
        }
    }
    std::cerr << "Sort Match: " << (match ? "OK" : "ERROR") << "\n";

    std::cerr << std::fixed << std::setprecision(3);
    std::cerr << "Bucket Sort: " << bucketTime.count() << " ms\n";
    std::cerr << "STL Stable Sort: " << stlTime.count() << " ms\n";

    for (int i = 0; i < realSize; ++i) {
        sorted[i].Print();
    }

    delete[] counts;
    delete[] offsets;
    delete[] currentOffset;
    delete[] sorted;
    delete[] raw;
    delete[] stlCopy;

    return 0;
}
\end{lstlisting}

\pagebreak

\section{Консоль}
\begin{alltt}
root@5703d59f664e:/workspaces/c_labs/da_lab1# g++ -O3 main.cpp -o main
root@5703d59f664e:/workspaces/c_labs/da_lab1# python3 test_gen.py 10000  
root@5703d59f664e:/workspaces/c_labs/da_lab1# ./main < test.txt > result.txt
Sort Match: OK
Bucket Sort: 1.011 ms
STL Stable Sort: 1.832 ms
root@5703d59f664e:/workspaces/c_labs/da_lab1# python3 test_gen.py 20000  
root@5703d59f664e:/workspaces/c_labs/da_lab1# ./main < test.txt > result.txt
Sort Match: OK
Bucket Sort: 1.957 ms
STL Stable Sort: 4.039 ms
root@5703d59f664e:/workspaces/c_labs/da_lab1# python3 test_gen.py 30000 
root@5703d59f664e:/workspaces/c_labs/da_lab1# ./main < test.txt > result.txt
Sort Match: OK
Bucket Sort: 3.117 ms
STL Stable Sort: 6.678 ms
root@5703d59f664e:/workspaces/c_labs/da_lab1# python3 test_gen.py 40000 
root@5703d59f664e:/workspaces/c_labs/da_lab1# ./main < test.txt > result.txt
Sort Match: OK
Bucket Sort: 4.127 ms
STL Stable Sort: 8.762 ms
root@5703d59f664e:/workspaces/c_labs/da_lab1# python3 test_gen.py 50000
root@5703d59f664e:/workspaces/c_labs/da_lab1# ./main < test.txt > result.txt
Sort Match: OK
Bucket Sort: 6.569 ms
STL Stable Sort: 12.071 ms
\end{alltt}
\pagebreak

\section{Тест производительности}

Тест производительности представляет собой измерение времени работы реализованного алгоритма и 
его сравнение со стандартным алгоритмом \texttt{std::stable\_sort} на наборах данных объемом от 10000 до 50000 пар \enquote{ключ-значение}. 
Для каждого теста генерировались случайные вещественные ключи в диапазоне $[-100, 100]$.
Ниже представлен график зависимости времени выполнения программы, на котором видно, что программа работает за линейное время, что соответствует теоретической оценке алгоритма карманной сортировки.

\begin{center}
\begin{tikzpicture}
    \begin{axis}[
        xlabel={Количество элементов},
        ylabel={Время (мс)},
        grid=major,
        legend pos=north west,
        width=12cm,
        height=7cm
    ]
    \addplot[color=blue, mark=square*] coordinates {
        (10000, 1.011)
        (20000, 1.957)
        (30000, 3.117)
        (40000, 4.127)
        (50000, 6.569)
    };
    \addlegendentry{Bucket Sort}

    \addplot[color=red, mark=x] coordinates {
        (10000, 1.832)
        (20000, 4.039)
        (30000, 6.678)
        (40000, 8.762)
        (50000, 12.071)
    };
    \addlegendentry{STL Stable Sort}
    \end{axis}
\end{tikzpicture}
\end{center}

\pagebreak

\section{Выводы}
Выполнив первую лабораторную работу по курсу \enquote{Дискретный анализ}, я изучил алгоритм карманной сортировки и его реализацию на языке C++. 
В процессе выполнения работы я на практике убедился, что если задано равномерное распределение входных ключей, данный алгоритм демонстрирует линейную сложность $O(n)$.

\pagebreak

\input{../src/literature}

\end{document}